%%%%%%%%%%%%%%%%%%%%%%%%%%%%%%%%%%%%%%%%%%%%%%%%%%%%%%%%%%%%
% (01) Input Settings
%%%%%%%%%%%%%%%%%%%%%%%%%%%%%%%%%%%%%%%%%%%%%%%%%%%%%%%%%%%%
\input{settings.tex}

%%%%%%%%%%%%%%%%%%%%%%%%%%%%%%%%%%%%%%%%%%%%%%%%%%%%%%%%%%%%
% (02) DOCUMENT INFORMATION - FROM YAML METADATA
%%%%%%%%%%%%%%%%%%%%%%%%%%%%%%%%%%%%%%%%%%%%%%%%%%%%%%%%%%%%
\newcommand{\DocTitle}{شی‌گرایی (چندریختی)}

\newcommand{\CourseTitle}{برنامه‌سازی پیشرفته و کارگاه}

\newcommand{\Semester}{بهار ۱۴۰۵}

\newcommand{\FirstProfessor}{دکتر احمدیان}

\newcommand{\SecondProfessor}{}

\newcommand{\Writers}{سام قربانی، آرمان حسینی، محمدحسین هاشمی، کیانا
پهلوان، کیانا رضوی، پیام علی‌مرادپور، مهدی جعفری}

\newcommand{\Faculty}{دانشکده ریاضی و علوم کامپیوتر}

\newcommand{\University}{دانشگاه صنعتی امیرکبیر}

%%%%%%%%%%%%%%%%%%%%%%%%%%%%%%%%%%%%%%%%%%%%%%%%%%%%%%%%%%%%
% DOCUMENT BODY
%%%%%%%%%%%%%%%%%%%%%%%%%%%%%%%%%%%%%%%%%%%%%%%%%%%%%%%%%%%%
\begin{document}
	
%%%%%%%%%%%%%%%%%%%%%%%%%%%%%%%%%%%%%%%%%%%%%%%%%%%%%%%%%%%%
% === TITLE PAGE ===
%%%%%%%%%%%%%%%%%%%%%%%%%%%%%%%%%%%%%%%%%%%%%%%%%%%%%%%%%%%%
\begin{titlepage}
	\centering
	\vspace*{1cm}
	
	%— Smaller logo (20% of text width)
	\begin{figure}[H]
		\centering
		
\includegraphics[width=0.20\textwidth]{images/logo-fa.png}
	\end{figure}
	
	\vspace{1cm}
	
	{\fontsize{18}{28}\selectfont \textbf{\CourseTitle} }
	
	\vspace{3cm}
	
	{\fontsize{22}{28}\selectfont \textbf{\DocTitle} }
	
	\vspace{3cm}
	
	{\fontsize{14}{22}\selectfont \rl{نگارش:}}\\
	\vspace{0.5cm}
	{\fontsize{16}{22}\selectfont \textbf{ \Writers } }

	
	\vspace{1cm}
	
	{\fontsize{14}{22}\selectfont \rl{استاد درس:}}\\
	\vspace{0.5cm}
	{\fontsize{16}{22}\selectfont \textbf{ \FirstProfessor } }
	
	
	\vfill
	
	{\fontsize{14}{18}\selectfont \textbf{\Faculty} }\\
	\vspace{0.5cm}
	{\fontsize{14}{18}\selectfont \textbf{\University} }
	\vspace{0.5cm}
	
	\par
	
	{\fontsize{14}{18}\selectfont \textbf{\Semester} }
	
\end{titlepage}

%— Switch to our fancy header/footer from this page onward
\clearpage

%%%%%%%%%%%%%%%%%%%%%%%%%%%%%%%%%%%%%%%%%%%%%%%%%%%%%%%%%%%%
% === TABLE OF CONTENTS ===
%%%%%%%%%%%%%%%%%%%%%%%%%%%%%%%%%%%%%%%%%%%%%%%%%%%%%%%%%%%%
\tableofcontents
\clearpage
\pagestyle{aftertitle}

%%%%%%%%%%%%%%%%%%%%%%%%%%%%%%%%%%%%%%%%%%%%%%%%%%%%%%%%%%%%
% === MAIN CONTENT ===
%%%%%%%%%%%%%%%%%%%%%%%%%%%%%%%%%%%%%%%%%%%%%%%%%%%%%%%%%%%%
\chapter{\DocTitle}
\section{مقدمه}

در این داکیومنت به مفاهیم پیشرفته‌تری از شیءگرایی می‌پردازیم. چندریختی
\lr{(Polymorphism)} یکی از این مفاهیم است که این امکان را می‌دهد که
متدهایی با نام‌های یکسان را به شکل‌های متفاوت پیاده‌سازی کنیم. اینترفیس نیز
به کلاس‌های متفاوت، این امکان را می‌دهد که متدهای مشترکی را بدون به اشتراک
گذاشتن کد و جزئیات پیاده‌سازی، داشته باشند. اینترفیس ها به ما کمک می‌کنند
برنامه ساختارمندتر و منعطف‌تری داشته باشیم.

حتماً کدهای این داک را اجرا کنید و هر سوالی که داشتید از تدریس‌یارها
بپرسید، در مورد آن سرچ کنید و از \lr{DeepSeek} راجع‌بهش اطلاعات بخواهید.

\section{\lr{Polymorphism}}

در دنیای واقعی وقتی اطرافتان را نگاه می‌کنید، می‌بینید که خیلی از کارهای
یکسان در موجودات مختلف به شکل‌های متفاوتی انجام می‌شود. مثلاً راه رفتن را
در نظر بگیرید؛ هم اسب، هم انسان هر دو راه می روند، اما میدانیم که در
انجام ان بسیار متفاوت‌اند یا مثلاً غذا خوردن، حرف زدن و خصوصیات دیگری که
در موجودات مشترک‌اند اما در جزئیات و عملکرد متفاوت‌اند.

حال در دنیای برنامه‌نویسی شی‌گرا، اگه بخواهید یک متد، کلاس یا شیء در شرایط
مختلف رفتارهای متفاوتی داشته باشد، از قابلیتی به اسم \lr{Polymorphism}
استفاده می‌کنید. این ویژگی باعث می‌شود کد شما منعطف‌تر، خواناتر و نگهداری
آن راحت‌تر شود.

جاوا دارای دو نوع پلی مورفیسم است:

\begin{itemize}
\item
  \textbf{\lr{Compile-time Polymorphism (Method Overloading)}}:

  در این حالت چند متد با یک اسم ولی با ورودی‌های متفاوت در یک کلاس تعریف
  می‌کنید. این کار به شما اجازه می‌دهد که یک کار مشابه را با ورودی‌های
  مختلف انجام بدهید.
\item
  \textbf{\lr{Run-time Polymorphism (Method Overriding)}}:
\end{itemize}

این یعنی یک متد در کلاس والد تعریف شده است و بعد در کلاس فرزند با همان
اسم ولی پیاده‌سازی متفاوت بازنویسی می‌شود. این کار به شما اجازه می‌ده رفتار
یه متد را در کلاس‌های فرزند تغییر بدهید.

ایده اصلی این است که متدهایی با اسم یکسان می‌توانند رفتارهای متفاوتی
داشته باشند و همین باعث می‌شود کد شما هم تمیزتر و هم راحت‌تر قابل استفاده
مجدد باشد.

\subsection{\lr{Compile-time Polymorphism (Method Overloading)}}

گاهی وقت ها پیش میاید که چند متد با یک اسم در یک کلاس تعریف می‌کنیم، اما
امضای (تعداد یا نوع پارامترهای) متدها با هم متفاوت اند. این کار باعث
می‌شود وقتی متد را صدا می‌زنیم، کامپایلر بفهمد کدام متد باید اجرا شود. به
این روش \lr{Method Overloading} می‌گویند.

ویژگی‌های این نوع polymorphism، به شکل زیر است:

\begin{itemize}
\tightlist
\item
  در یک کلاس انجام می‌شود.
\item
  تفاوت در تعداد یا نوع پارامترها وجود دارد.
\item
  در زمان کامپایل مشخص می‌شود.
\end{itemize}

مثلا در کد زیر، چند متد \lr{\texttt{add}} داریم که هر کدام، ورودی‌های
متفاوتی می‌گیرند و خروجی متفاوتی می‌دهند:

\begin{codebox}{java}
class MathOperations {  
    int add(int a, int b) {  
        return a + b;  
    }  

    double add(double a, double b) {  
        return a + b;  
    }  

    int add(int a, int b, int c) {  
        return a + b + c;  
    }  
}

public class Main {  
    public static void main(String[] args) {  
        MathOperations math = new MathOperations();  
        System.out.println(math.add(5, 10)); // Output: 15  
        System.out.println(math.add(5.5, 2.2)); // Output: 7.7  
        System.out.println(math.add(1, 2, 3)); // Output: 6  
    }
}
\end{codebox}

\subsection{\lr{Run-time Polymorphism (Method Overriding)}}

اینجا ایده این است که یک متد توی کلاس \lr{Parent} داریم و یه کلاس
\lr{Child} که می‌خواهد آن متد را به شکل خودش پیاده‌سازی کند. این کار را با
\lr{Method Overriding} انجام می‌دهیم.

ویژگی‌های این نوع polymorphism، به شکل زیر است:

\begin{itemize}
\tightlist
\item
  بین کلاس \lr{Parent} و \lr{Child} اتفاق می‌افتد.
\item
  امضای متد (شامل اسم آن، جنس ورودی‌ها و خروجی‌ها) تغییری نمی‌کند.
\item
  در زمان اجرای برنامه‌، بسته به تایپ متغیر تصمیم گرفته می‌شود کدام متد
  اجرا شود.
\end{itemize}

مثلا در کد زیر، متد \lr{\texttt{makeSound}} در \lr{\texttt{Dog}} دوباره
تعریف شده است. لازم است که وقتی متدی را \lr{override} می‌کنید، حتما بالای
آن از \lr{\texttt{@Override}} استفاده کنید، وگرنه \lr{IntelliJ} به شما
هشدار می‌دهد.

\begin{codebox}{java}
class Animal {
    void makeSound() {  
        System.out.println("An animal makes a sound");
    }  
}  
  
class Dog extends Animal {
    @Override
    void makeSound() {  
        System.out.println("Dog barks");
    }  
}  
  
public class Main {
    public static void main(String[] args) {  
        Animal myAnimal = new Animal();
        myAnimal.makeSound();  // Output: An animal makes a sound
        
        Animal myDog = new Dog();
        myDog.makeSound();  // Output: Dog barks    }  
}
\end{codebox}

\section{\lr{interface}}

ارث‌ بری و وراثت ابزار بسیار کاربردیه که به شما اجازه میدهد از
\lr{polymorphism} (چند ریختی) استفاده کنید ولی خب اگر توجه کرده باشید،
محدودیتی هم جلوی شما قرار می دهد. جاوا از ارث بری یگانه
\lr{(single inheritance)} استفاده میکند، به این معنا که یک کلاس میتواند
فقط یک \lr{superclass} را \lr{extend} کند. این محدودیت، اجازه نمیدهد که
از ارث بری برای کلاس‌هایی استفاده کنیم که نیاز دارند شاخصه‌های متفاوت
موجود در سوپر کلاس‌های مختلف را همزمان داشته باشند. امتحان کنید و ببینید
که اگه بخواهید همزمان چند تا \lr{superclass} را \lr{extend} کنید، به
کامپایل ارور بر میخورید.

همچنین، در بعضی موقعیت‌ها، ممکن است نخواهید که جزییات پیاده‌سازی و کد را
به اشتراک بگذارید. اینجاست که با مفهوم جدیدی به اسم \lr{interface} آشنا
میشویم تا محدودیت‌های قبلی را کنار بگذاریم.

\subsection{اینترفیس چیست و چطور از آن استفاده کنیم؟}

اینترفیس مثل کلاس‌های عادی است، صرفا یکسری تفاوت‌ها و محدودیت‌ها دارد.
\lr{Interface} متشکل شده از یکسری تعریف اولیه از متدها است. متدهایی که
بدنه ندارند و جزییات پیاده سازی آنها نوشته نشده است.

کلاس‌ها، سوپرکلاس‌های عادی یا \lr{abstract} را \lr{extend} میکردند و
ویژگی‌های آنها را به ارث میبردند، ولی \lr{interface} ها را \lr{implement}
(پیاده‌سازی) میکنند. \lr{Implement} کردن یک اینترفیس به این معناست که
کلاسی که میخواهد \lr{interface} مورد نظر را پیاده‌سازی کند، قول میدهد که
بدنه همه متدهای تعریف شده در اینترفیس را پیاده‌سازی کند (در غیر این صورت
به کامپایل ارور میخورد).

به عنوان مثال، میخواهیم یک اینترفیس به نام \lr{\texttt{shape}} بسازیم.
تصور کنید که میخواهیم کلاس‌هایی را بسازیم که هر کدام نماینده یک شکل
متفاوت باشند (مثلا مستطیل و دایره). ارث‌بری ایده خوبیست چون که هر دوی این
اشکال رفتار مشابهی دارند (مثل مساحت و محیط) ولی نحوه محاسبه این مساحت و
محیط (پیاده‌سازی متدها)، متفاوت و مرتبط با نوع شکل است.به این شکل، سلسه
مراتب پولی‌مورفیسمی از کلاس‌های \lr{\texttt{shape}} داریم بدون اینکه کدی و
پیاده‌سازی‌ای بین آن ها به اشتراک گذاشته شده باشد.

برای نوشتن یک اینترفیس، از کلمه کلیدی \lr{\texttt{interface}} استفاده
میکنیم:

\begin{codebox}{java}
public interface Shape {  
}
\end{codebox}

در داخل بدنه interface، تعریف متدها را مینویسیم ولی بجای اینکه از
\lr{\texttt{\{\}}} استفاده کنیم و داخل \lr{braces} پیاده‌سازی متد را
بنویسیم، به یک \lr{\texttt{;}} اکتفا میکنیم:

\begin{codebox}{java}
public interface Shape {  
    public double getArea();  
    public double getPerimeter();  
}
\end{codebox}

خب ما الان اولین اینترفیسمون را ساختیم. توجه کنید که \textbf{هم اینترفیس
هم \lr{class abstract} قابلیت این رو ندارند که از آنها آبجکت ساخته شود}.
علت این موضوع، وجود \lr{abstract methods} در هر دوی آنها است. پس اول
باید یک کلاس، کلاس \lr{abstract} را به ارث ببرد و یا اینترفیس را پیاده
سازی کند تا بعد به صورت غیر مستقیم، بتواند از کلاس ثانویه آبجکت ایجاد
کند.

حالا میخواهیم دو کلاس \lr{\texttt{Rectangle}} و \lr{\texttt{Circle}} را
بسازیم که اینترفیس \lr{\texttt{shape}} را \lr{implement} کنند. هنگام
تعریف این کلاس‌ها، باید از کلمه کلیدی \lr{implements} استفاده کنیم. دقت
کنید که این دو کلاس را باید در دو فایل جدا تعریف کنید:

\begin{codebox}{java}
public class Rectangle implements Shape {

}

  
public class Circle implements Shape {  

}
\end{codebox}

وقتی که در تعریف این دو کلاس، اینترفیس \lr{\texttt{shape}} را
\lr{implement} کردیم، تضمین کردیم که متدهای \lr{\texttt{getArea()}} و
\lr{\texttt{getPerimeter()}} را پیاده سازی کنیم و تا زمانی که این کار را
نکنیم، با کامپایل ارور روبه‌رو میشویم. کلاس‌ها را به صورت زیر گسترش
میدهیم:

\begin{codebox}{java}
public class Rectangle implements Shape{
    private double width;
    private double height;
    public Rectangle(double width, double height){ 
        this.width = width;
        this.height = height;  
    } 

    @Override
    public double getArea() {
        return width * height;  
    }

    @Override
    public double getPerimeter() {
        return 2.0 * (height + width);  
    }  
}

public class Circle implements Shape{
    private double radius;

    public Circle(double radius){
        this.radius = radius;  
    }

    @Override
    public double getArea() {
        return Math.PI * radius * radius;  
    }

    @Override
    public double getPerimeter() {
        return 2.0 * Math.PI * radius;  
    }  
}
\end{codebox}

الان کلاس‌های ما کامل شدند و میتوانیم ازشون استفاده کنیم و از روی آن ها
آبجکت بسازیم. به مثال زیر دقت کنید:

\begin{codebox}{java}
public class ShapeMain {
    public static void main(String[] args){  
        Shape circle = new Circle(5);  
        Shape rectangle = new Rectangle(3, 6);  
        System.out.println("rectangle's area = " + rectangle.getArea());  
        System.out.println("rectangle's perimeter = " + rectangle.getPerimeter());  
        System.out.println("circle's area = " + circle.getArea());  
        System.out.println("circle's perimeter = " + circle.getPerimeter());  
    }  
}
\end{codebox}

خروجی کد بالا:

\begin{codebox}{text}
rectangle's area = 18.0
rectangle's perimeter = 18.0
circle's area = 78.53981633974483
circle's perimeter = 31.41592653589793
\end{codebox}

در این مثال، یک آبجکت از کلاس \lr{\texttt{Rectangle}} و یک آبجکت از کلاس
\lr{\texttt{Circle}} ساختیم که هر دو اینترفیس \lr{\texttt{Shape}} را
\lr{implement} میکردند و پیاده سازی مخصوص به خود برای متد های این
اینترفیس را داشتن. به سینتکس دقت کنید.

کلاس‌هایی که اینترفیس یکسانی را \lr{implement} میکنند، سلسله مراتبی شبیه
به کلاس‌های فرزندی که یک کلاس پدر را به ارث میبرند، دارند. مثلا برای مثال
بالا، میتوانیم چنین ارتباطی را در نظر بگیریم:

\begin{figure}[H]
\centering
\includegraphics[width=0.8\textwidth]{images/oop-2-10.png}
\end{figure}

حال که میدانیم \lr{interface} چیست. به بررسی کاربرد آن و اینکه از آن در
چه جاهایی استفاده می شود میرسیم.

\subsection{کاربردهای \lr{interface}}

\begin{enumerate}
\def\labelenumi{\arabic{enumi}.}
\tightlist
\item
  اینترفیس مثل یک قرارداد یا پروتکل عمل می‌کند. در واقع یک ساختار و
  پروتوتایپ ارائه می‌دهد که کلاس‌های مرتبط با آن باید آن‌ها را داشته باشند.
  با این حال، هر کدام می‌توانند پیاده‌سازی مخصوص خودشان را ارائه دهند. به
  همین دلیل دسترسی بسیار مناسبی برای \lr{polymorphism} فراهم می‌کند، چون
  هر کلاسی که اینترفیس را پیاده‌سازی می‌کند می‌تواند پیاده‌سازی منحصر‌به‌فرد
  خودش را برای متدهای آن اینترفیس داشته باشد. همچنین این ویژگی امکان
  انعطاف‌پذیری بیشتری را در طراحی سیستم‌ها فراهم می‌کند.
\item
  با تعریف یک interface، جزئیات پیاده‌سازی و کد ارائه نمی‌شود و صرفاً یک
  قالب در اختیار قرار می‌گیرد.
\item
  مهم‌تر از همه، \textbf{چندگانه بودن وراثت \lr{(multiple inheritance)}}
  در استفاده از اینترفیس‌ها است. اینترفیس‌ها این امکان را فراهم می‌کنند که
  یک کلاس بتواند همزمان از چندین اینترفیس پیروی کرده و آن‌ها را
  \lr{implement} کند، اما در عین حال فقط می‌تواند یک کلاس را \lr{extend}
  کند.
\item
  همچنین استفاده از اینترفیس‌ها باعث کاهش وابستگی بین کلاس‌ها می‌شود، زیرا
  دیگر کلاس‌ها به جزئیات پیاده‌سازی یکدیگر وابسته نیستند. این موضوع
  می‌تواند منجر به کدهای تمیزتر و قابل نگهداری تر شود.
\item
  اینترفیس‌ها معمولاً برای جداسازی مسئولیت‌ها استفاده می‌شوند. به جای اینکه
  یک کلاس تمام رفتارها و ویژگی‌ها را در خود داشته باشد، اینترفیس‌ها امکان
  تقسیم مسئولیت‌ها را فراهم می‌کنند و این موضوع باعث می‌شود کدها قابل تست‌تر
  و نگهداری آن‌ها آسان‌تر شود.
\end{enumerate}

حالا یک سوالی که پیش می آید این است که خب چرا اصلاً \lr{Abstract Class}
وجود دارد وقتی که اینترفیس می‌تواند تقریباً همان کار را انجام بدهد؟ چرا
فقط یکی از این دو را نداریم؟ برای جواب دادن به این سوال، باید تفاوت‌ها را
بررسی کنیم. چون که این دو، برای اهداف متفاوتی طراحی شدند و نمیتوانند
کاملا جایگزین همدیگه بشوند.

\section{مقایسه \lr{interfaces} و \lr{abstract classes}}

کلاس \lr{abstract} برای ساختارهای سلسله‌مراتبی \lr{(Hierarchy)} که
اشتراکات زیادی دارند، مناسب است. در مقابل، اینترفیس برای قراردادهای
عمومی بین کلاس‌های نامرتبط استفاده می‌شود.

\begin{table}[H]
	\centering
	\small 
	\setLTR
	\begin{tabularx}{\textwidth}{>{\raggedright\arraybackslash}X>{\raggedright\arraybackslash}X}
		\hline
		\textbf{\lr{Interface}} & \textbf{\lr{Abstract Class}} \\
		\hline
		\lr{Interface can have only \textbf{abstract methods}.}
		&
		\lr{Abstract class can have \textbf{abstract and non-abstract methods}.} \\
		\hline
		\lr{Interface \textbf{supports multiple inheritance}.}
		&
		\lr{Abstract class \textbf{doesn't support multiple inheritance}.} \\
		\hline
		\lr{Interface has \textbf{only static and final variables}.}
		&
		\lr{Abstract class can have \textbf{final, non-final, static, and non-static variables}.} \\
		\hline
	\end{tabularx}
\end{table}

همونطور که در جدول بالا میبینید، اینترفیس‌ها و کلاس‌های \lr{abstract}
تفاوت‌های ملموسی دارند.

به جز تفاوت در نحوه وراثت و ارث‌بری، اینترفیس‌ها برخلاف کلاس‌های abstract،
فقط میتوانند فیلدهای \lr{static} و \lr{final} داشته باشند. یعنی
اینترفیس‌ها، متغیر نمونه \lr{(instance variable)} ندارند و فقط درصورتی
میتوانند فیلد داشته باشند که مقدارش ثابت و \lr{static} باشد.

یک فیلد مثل سن در یک اینترفیس‌ فقط میتواند به قالب زیر تعریف بشود:

\begin{codebox}{java}
public static final int age = 20;
\end{codebox}

این محدودیت به دلایل زیر به وجود آمده است:

\begin{itemize}
\tightlist
\item
  \textbf{اینترفیس‌ها برای «تعریف رفتار» ساخته شدند، نه ``نگه‌داری
  وضعیت'':}
\end{itemize}

اینترفیس‌ها یک قرارداد \lr{(Contract)} برای کلاس‌های پیاده‌سازشون ارائه
می‌کنند و نباید شامل وضعیت \lr{(State)} یا داده‌های متغیر باشند. در
اینترفیس‌ها، تمرکز بر روی متدهاست، نه متغیرهای نمونه . اگر اینترفیس‌ها
اجازه‌ی داشتن متغیرهای نمونه رو داشتند، در عمل شبیه به کلاس‌های
\lr{abstract} می‌شدند که خلاف فلسفه‌ی اینترفیس است.

\begin{itemize}
\tightlist
\item
  \textbf{اینترفیس‌ها نباید وابستگی به مقدارهای متغیر داشته باشند:}
\end{itemize}

اگر اینترفیس‌ها میتوانستند متغیر نمونه داشته باشند، هر کلاس پیاده‌سازی
آنها باید این مقدار رو مدیریت می‌کرد که باعث نقض استقلال اینترفیس می‌شد.

\begin{itemize}
\tightlist
\item
  \textbf{جلوگیری از مشکل «وراثت چندگانه» و «داده‌ی متناقض»:}
\end{itemize}

اگر اینترفیس‌ها متغیرهای نمونه داشتند، مشکل تناقض داده‌ها به وجود می‌آمد.
مثالی رو در نظر بگیرید که در آن، یه کلاس دو اینترفیس را پیاده‌سازی میکند
که هر کدام دارای یه فیلد با نام یکسان هستند ولی مقادیر متفاوتی دارند.

یک نکته قابل توجه دیگر، این است که اینترفیس بر خلاف \lr{class abstract}
کانستراکتور ندارد.

\textbf{حال که میدانیم این دو متفاوت اند و اینکه علت تفاوت چیست،
میخواهیم ببینیم که کجا و چرا باید از کدام ساختار استفاده کنیم.}

\begin{table}[H]
	\centering
	\small
	\begin{tabularx}{\textwidth}{>{\raggedright\arraybackslash}X>{\raggedright\arraybackslash}X}
		\hline
		\textbf{\lr{Abstract class}} & \textbf{\lr{Interface}} \\
		\hline
		اگر فقط به یک لایه وراثت نیاز دارید و مشخص است که کلاس‌های فرزند شباهت‌های زیادی در پیاده‌سازی بین خودشون دارن، یک کلاس ابسترکت می‌تواند مناسب‌تر باشد. &
		اگر نیاز دارید یک کلاس از چند منبع مختلف ویژگی‌هایی دریافت کند، اینترفیس‌ها گزینه بهتری هستند. \\
		\hline
		اگر برخی متدها باید در همه‌ی کلاس‌های فرزند یکسان باشند و نمی‌خواهید کد تکراری بنویسید، کلاس ابسترکت بهتر است. &
		اگر فقط می‌خواهید یک قرارداد تعیین کنید که همه‌ی کلاس‌ها موظف باشند متدهایی را پیاده‌سازی کنند و لزوما این پیاده‌سازی‌ها شبیه هم نیستن، اینترفیس گزینه‌ی بهتری است. \\
		\hline
		اگر کلاس باید مقادیر پیش‌فرض یا داده‌ای را در خود نگه دارد، کلاس ابسترکت بهتر است. مزیت کلاس‌های ابسترکت این است که میتوانیم هم متغیرهای مشترک و هم متدهای مشترک با سوپرکلاس را تعریف کنیم و در صورت نیاز به تغییر پیاده‌سازی در یک کلاس فرزند، \lr{override} کنیم. &
		اگر فقط متدهای مورد نیاز را تعریف می‌کنید و هیچ داده‌ و فیلدی در کلاس پایه نیاز ندارید (نیاز به متغیرهای نمونه ندارید) اینترفیس بهترین انتخاب است. \\
		\hline
	\end{tabularx}
\end{table}

\subsection{\lr{Multiple inheritance by Interfaces}}

حالا میخواهیم یک مثال را بررسی کنیم که یک کلاس، چند تا اینترفیس مختلف را
\lr{implement} میکند:

تصور کنید یک شرکت فناوری قصد دارد یک ربات چندمنظوره بسازد که بتواند:

\begin{enumerate}
\def\labelenumi{\arabic{enumi}.}
\tightlist
\item
  پرواز کند (مثل یک پهپاد)
\item
  شنا کند (مثل یک زیردریایی)
\item
  صحبت کند و ارتباط برقرار کند (مانند یک دستیار هوشمند)
\end{enumerate}

هر یک از این قابلیت‌ها شامل چندین رفتار است. ما این قابلیت‌ها را به‌عنوان
اینترفیس‌ها تعریف می‌کنیم تا ربات بتواند از اونا به‌صورت ارث‌بری چندگانه
استفاده کند:

\begin{codebox}{java}
// interface for flying feature
interface Flyable {
    void fly();         // start flyting
    void adjustAltitude(int height); // adjusting height if flying
    void land();        // start landing
}

// interface for swim feature
interface Swimmable {
    void swim();        // start swimming
    void adjustDepth(int depth); // adjusting depth of swim
    void stopSwimming(); // stopping the swimming process
}

// interface for speaking feature
interface Speakable {
    void speak(String message); // speaking
    void setLanguage(String language); // changing the language
    void processCommand(String command); // processing voice commands
}

package Robot;
import java.lang.String;

// The class which implements 3 interfaces
class SmartRobot implements Flyable, Swimmable, Speakable {

    private int altitude;
    private int depth;
    private String language;
    public SmartRobot(int altitude, int depth, String language){
        this.depth = depth;
        this.altitude = altitude;
        this.language = language;  
    }

    //implementing fly methods
    @Override
    public void fly() {  
        System.out.println("Robot is flying ...");  
    }

    @Override
    public void adjustAltitude(int height) {
        this.altitude = height;  
        System.out.println("Height has been set to " + height + " meters.");  
    }

    @Override
    public void land() {
        this.altitude = 0;  
        System.out.println("Robot has landed successfully");  
    }

    // implementing swim methods
    @Override
    public void swim() {  
        System.out.println("Robot is swimming");  
    }

    @Override
    public void adjustDepth(int depth) {
        this.depth = depth;  
        System.out.println("Depth has been set to" + depth + " meters");  
    }

    @Override
    public void stopSwimming() {  
        System.out.println("Robot stopped swimming successfully");  
    }

    // implementing speak methods
    @Override
    public void speak(String message) {  
        System.out.println("Robot says: " + message);  
    }

    @Override
    public void setLanguage(String language) {
        this.language = language;  
        System.out.println("Language has been change into " + language);  
    }

    @Override
    public void processCommand(String command) {  
        System.out.println("Processing the command: " + command);  
    }  
}
\end{codebox}

تو این مثال من همه بخش هارو پشت هم نوشتم ولی در اصل بهتره هر اینترفیس و
کلاس رو تو یه فایل جداگانه تعریف کنید.

\subsection{بعضی اینترفیس‌های مهم در جاوا}

خیلی از کلاس‌های موجود در لایبرری‌های جاوا، همزمان یک سوپرکلاس را
\lr{extend} کرده و یک یا بعضا چند اینترفیس را \lr{implement} میکنند.
برای مثال می توان به کلاس \lr{\texttt{PrintStream}} اشاره کرد ( کلاسی که
\lr{\texttt{system.out}} یک نمونه \lr{(instance)} از آن است). این کلاس
به این صورت تعریف می شود:

\begin{codebox}{java}
public class PrintStream extends FilterOutputStream implements Appendable, closeable {

}
\end{codebox}

یک دسته اینترفیس مهم دیگه هم که جاوا در اختیارتون میگذاره، اینترفیس‌های
\lr{\texttt{List}}، \lr{\texttt{Map}} و \lr{\texttt{set}} هستند که توسط
پکیج \lr{\texttt{java.util}} ارائه می شوند و با \lr{implement} کردن آنها
میتوانید ساختمان داده‌هایی را توصیف کنید که به کمک آنها مجموعه‌ای از
آبجکت‌ها را ذخیره کنید.

\section{\lr{Casting} در جاوا}

قبل از ادامهٔ ارث‌بری، مقداری در مورد \lr{cast} کردن صحبت می‌کنیم. اصطلاح
\lr{cast} کردن، برای تغییر تایپ متغیرها استفاده می‌شود. مثلا این کد را
ببینید:

\begin{codebox}{java}
public static void main(String[] args) {
    int a = 10;
    long b = a;  
  
    System.out.println("a: " + a);  
    System.out.println("b: " + b);  
}
\end{codebox}

در این کد، متغیر \lr{\texttt{a}} با تایپ \lr{\texttt{int}} به متغیری از
جنس \lr{\texttt{long}} به اسم \lr{\texttt{b}} تبدیل شده است. این تبدیل
بدون هیچ مشکلی انجام شده چون بازهٔ اعدادی که متغیرهای \lr{\texttt{long}}
پوشش می‌دهند بسیار بزرگتر از اعداد \lr{\texttt{int}} است. به خاطر درستی
این cast، خروجی این کد هم به شکل زیر است:

\begin{codebox}{text}
a: 10  
b: 10
\end{codebox}

حالا اگر متغیر \lr{\texttt{long}} را به یک متغیر \lr{\texttt{int}}، کست
کنیم چه اتفاقی می‌افتد؟ کد زیر را ببینیند:

\begin{codebox}{java}
public static void main(String[] args) {
    long a = 10_000_000_000L;
    int b = a;  
  
    System.out.println("a: " + a);  
    System.out.println("b: " + b);  
}
\end{codebox}

تلاش کنید این کد را اجرا کنید. به خطای زیر می‌خورید:

\begin{figure}[H]
\centering
\includegraphics[width=0.8\textwidth]{images/oop-2-1.png}
\end{figure}

متن خطا می‌گوید:

\begin{codebox}{text}
java: incompatible types: possible lossy conversion from long to int
\end{codebox}

جاوا با این خطا به شما می‌گوید که تبدیل متغیرهای \lr{\texttt{long}} به
\lr{\texttt{int}}، خطرناک است و می‌تواند منجر به خطاهای \lr{overflow}
شود. مثلاً در همین مثال، عدد \lr{\texttt{10000000000}} به قدری بزرگ است
که در یک متغیر \lr{\texttt{int}} جا نمی‌شود و به همین خاطر تبدیل آن به یک
\lr{\texttt{int}} مشکلاتی برای شما به همراه دارد.

حالا اگر یک \lr{\texttt{int}} قبل از \lr{\texttt{a}} قرار دهیم به جاوا
می‌گوییم که از \lr{cast} کردن مطمئنیم و می‌دانیم که هیچ مشکلی ایجاد
نمی‌کند:

\begin{codebox}{java}
public static void main(String[] args) {
    long a = 10_000_000_000L;
    int b = (int)a;  
  
    System.out.println("a: " + a);  
    System.out.println("b: " + b);  
}
\end{codebox}

خروجی این کد به شکل زیر است:

\begin{codebox}{text}
a: 10000000000  
b: 1410065408
\end{codebox}

حالا به مقدار عجیب \lr{\texttt{b}} نگاه کنید! مشکلی که می‌بینید، مشکل
بسیار معروفی به اسم ``\lr{overflow}'' یا ``سرریزی'' است. زمانی که شما،
عدد خیلی بزرگی را داخل یک متغیر که فضای کافی ندارد ذخیره می‌کنید، به خطای
\lr{overflow} می‌خورید. خوب است بدانید مقدار چاپ شدهٔ
\lr{\texttt{1410065408}} دقیقاً برابر \lr{\texttt{31}} بیت سمت راست
\lr{\texttt{10000000000}} می‌باشد. (یادآوری: \lr{\texttt{int}} معمولاً
\lr{\texttt{32}} بیت در حافظه اشغال می‌کند؛ \lr{\texttt{1}} بیت علامت و
\lr{\texttt{31}} بیت برای مقدار)

حالا اگر مقدار \lr{\texttt{a}} کوچکتر بود به این خطا نمی‌خوردید:

\begin{codebox}{java}
public static void main(String[] args) {
    long a = 100;
    int b = (int)a;  
  
    System.out.println("a: " + a);  
    System.out.println("b: " + b);  
}
\end{codebox}

خروجی این کد به شکل زیر است:

\begin{codebox}{text}
a: 100  
b: 100
\end{codebox}

علت این که \lr{cast} بالا به درستی انجام شد این است که مقدار
\lr{\texttt{100}} به راحتی در یک متغیر \lr{\texttt{int}} جا می‌شود.

اتفاق مشابهی در \lr{cast} کردن بین متغیرهای \lr{\texttt{int}} و
\lr{\texttt{double}} هم رخ می‌دهد. \lr{Cast} زیر به درستی انجام می‌شود:

\begin{codebox}{java}
public static void main(String[] args) {
    int a = 100;
    double b = a;  
  
    System.out.println("a: " + a);  
    System.out.println("b: " + b);  
}
\end{codebox}

خروجی این کد به صورت زیر است:

\begin{codebox}{text}
a: 100  
b: 100.0
\end{codebox}

ممیز خط دوم نشان می‌دهد که متغیر \lr{\texttt{b}} از جنس
\lr{\texttt{double}} است.

ولی اگر تلاش کنید عددی اعشاری را به عددی صحیح تبدیل کنید، باز هم خطا
می‌گیرید:

\begin{codebox}{java}
public static void main(String[] args) {
    double a = 3.14;
    int b = a;  
  
    System.out.println("a: " + a);  
    System.out.println("b: " + b);  
}
\end{codebox}

دوباره با همان خطای قبلی روبه‌رو شدیم:

\begin{codebox}{text}
java: incompatible types: possible lossy conversion from double to int
\end{codebox}

این خطا به خاطر این است که در این cast، بخش اعشاری متغیر \lr{\texttt{a}}
از دست می‌رود. اگر شما این \lr{cast} را با استفاده از \lr{\texttt{int}}
به صورت زیر انجام دهید:

\begin{codebox}{java}
int b = (int)a;
\end{codebox}

می‌بینید که فقط بخش صحیحِ متغیر \lr{\texttt{a}}، در \lr{\texttt{b}} ذخیره
می‌شود:

\begin{codebox}{text}
a: 3.14  
b: 3
\end{codebox}

حتی اگر دوباره متغیر \lr{\texttt{b}} را به \lr{\texttt{double}} تبدیل
کنید، می‌بینید که بخش اعشاری‌ برنمی‌گردد:

\begin{codebox}{java}
public static void main(String[] args) {
    double a = 3.14;
    int b = (int)a;
    double c = b;  
  
    System.out.println("a: " + a);  
    System.out.println("b: " + b);  
    System.out.println("c: " + c);  
}
\end{codebox}

خروجی این کد به شکل زیر است:

\begin{codebox}{text}
a: 3.14  
b: 3  
c: 3.0
\end{codebox}

این cast، خطرناک محسوب می‌شود چون در اون، بخشی از داده‌ها از دست رفته‌اند.

\subsection{انواع \lr{casting}}

دیدید که \lr{cast} کردن، می‌تواند خطرناک باشد. جاوا می‌داند که چه
\lr{cast} هایی می‌توانند مشکل‌زا باشند. اگر مثل مثال اول، یک متغیر
\lr{\texttt{int}} را به \lr{\texttt{long}} تبدیل کنیم، جاوا با اطمینان
از موفق و بی‌خطر بودن این \lr{cast} ، آن را بدون هیچ خطا و اخطاری اجرا
می‌کند. اما در مثال دوم، که \lr{cast} کردن می‌توانست مشکل‌زا باشد، جاوا به
ما خطا داد و مجبور شدیم با استفاده از \lr{\texttt{(int)}} به او بگوییم
که ما نسبت به این مشکلات آگاهیم و همچنان می‌خواهیم این \lr{cast} را انجام
دهیم. به طور کلی \lr{cast} هایی که ممکن است برای برنامهٔ شما با از دست
دادن اطلاعات، \lr{overflow} یا هر خطای دیگری مواجه شوند، \lr{cast} های
خطرناکی محسوب می‌شوند.

بر همین اساس، \lr{cast} ها را به دو دسته تقسیم می‌کنیم:

\subsubsection{کست‌های \lr{implicit}}

واژهٔ ``\lr{implicit}'' یعنی چیزی که بدون اشارهٔ مستقیم فهمیده می‌شود. مثل
وقتی که شما لازم نیست با گذاشتن \lr{\texttt{(int)}} مستقیما به جاوا بگید
که از خطرات \lr{cast} آگاهید. \lr{cast} هایی که مشکلی برای برنامهٔ شما
ایجاد نمی‌کنند، به \lr{cast} های \lr{implicit} معروف هستند. \lr{cast}
کردن متغیر \lr{\texttt{int}} به \lr{\texttt{long}} از همین نوع است.

\subsubsection{کست‌های \lr{explicit}}

واژهٔ ``\lr{explicit}'' به معنای چیزی است که باید مستقیماً و به شکل دقیق
بیان شود. کست‌های \lr{explicit} ، کست‌هایی هستند که ممکن است مشکل‌زا باشند
و لازم است که در کدتان حتما اعلام کنید که از انجام آن \lr{cast} آگاهید و
در جریان مشکلات احتمالی هستید. تبدیل کردن متغیر \lr{\texttt{long}} به
\lr{\texttt{int}} از این جنس است.

\subsection{کست کردن در \lr{reference type} ها}

حالا به ارث‌بری برگردیم. اگر به خاطر داشته باشید، به شما گفته بودیم
\lr{class} ها، ابزاری هستند که به وسیلهٔ آن‌ها می‌توانید برای خودتان
\lr{type} های جدید تعریف کنید. ولی خب، آیا این \lr{type} ها هم قابل
\lr{cast} کردن هستند؟ جواب طبیعتاً بله‌ست.

بیایید با یک مثال \lr{cast} کردن در این \lr{type} ها را ببینیم. کد زیر
را در تابع \lr{\texttt{main}} بنویسید:

\begin{codebox}{java}
public static void main(String[] args) {  
    Car myCar = new Car("Iran-Khodro", 180, 4);  
    myCar.showCarInfo();  
}
\end{codebox}

آن را اجرا کنید تا مطمئن باشید همه‌چیز درست کار می‌کند:

\begin{codebox}{text}
Brand: Iran-Khodro, Speed: 180 km/h  
Doors: 4
\end{codebox}

حالا، بیاید یک \lr{cast} عجیب بین \lr{\texttt{Car}} و
\lr{\texttt{Vehicle}} انجام بدیم:

\begin{codebox}{java}
public static void main(String[] args) {  
    Car myCar = new Car("Iran-Khodro", 180, 4);  
    myCar.showCarInfo();  
  
    System.out.println();  
  
    Vehicle myVehicle = myCar; // Casting from Car to Vehicle
    myVehicle.showInfo();  
}
\end{codebox}

خروجی این کد، به شکل زیر است:

\begin{codebox}{text}
Brand: Iran-Khodro, Speed: 180 km/h  
Doors: 4  

Brand: Iran-Khodro, Speed: 180 km/h
\end{codebox}

اگر دقت کنید، وقتی که \lr{\texttt{Car}} از \lr{\texttt{Vehicle}} ارث‌بری
می‌کند، تمام فیلدها و متدهای \lr{\texttt{Vehicle}} متعلق به
\lr{\texttt{Car}} نیز هستند. یعنی هر آبجکتی از جنس \lr{\texttt{Car}} ،
تمام فیلدهای یک \lr{\texttt{Vehicle}} را در خودش دارد، پس ما همه چیز
برای \lr{cast} کردن \lr{\texttt{Car}} به \lr{\texttt{Vehicle}} را داریم
و می‌توانیم آن را انجام دهیم. همان‌طور که می‌بینید، برند
\lr{\texttt{myVehicle}} ایران خودرو و سرعتش \lr{\texttt{180 km/h}} است.
دقیقا مثل \lr{\texttt{myCar}} .

حتی می‌توانیم \lr{\texttt{myVehicle}} را با یک دوچرخه هم جایگزین کنیم:

\begin{codebox}{java}
public static void main(String[] args) {  
    Car myCar = new Car("Iran-Khodro", 180, 4);  
    myCar.showCarInfo();  
  
    System.out.println();  
  
    Vehicle myVehicle = myCar; // Casting from Car to Vehicle
    myVehicle.showInfo();  
  
    System.out.println();  
  
    Bike myBike = new Bike("Scott", 10, false);  
    myBike.showBikeInfo();  
  
    System.out.println();  
  
    myVehicle = myBike; // Casting, again!
    myVehicle.showInfo();  
}
\end{codebox}

می‌بینید که \lr{\texttt{myVehicle}} از یک ماشین به یک دوچرخه تبدیل می‌شود.

بیاید یک لحظه به اتفاق خیلی خوبی که الآن دیدیم بیشتر فکر کنیم. دیدیم که
متغیرهایی با جنس \lr{\texttt{Vehicle}} ، می‌توانند با آبجکت‌های
\lr{\texttt{Car}} ، \lr{\texttt{Bike}} و \lr{\texttt{MotorCycle}} پر
شوند. اتفاقی که بدون ارث‌بری ممکن نبود. فرض کنید بدون ارث‌بری، می‌خواستید
متدی بنویسید که با توجه به سرعت یک وسیلهٔ نقلیه، بگوید یک مسافت مشخص در
چند ثانیه طی می‌شود. اگر ارث‌بری نبود، مجبور بودید به شکل زیر، برای هر یک
از وسایل نقلیه‌ متدی جدا تعریف کنید:

\begin{codebox}{java}
public static double calculateTimeForCar(double distance, Car car) {
    return distance / car.speed;  
}

public static double calculateTimeForBike(double distance, Bike bike) {
    return distance / bike.speed;  
}

public static double calculateTimeForMotorCycle(double distance, MotorCycle motorCycle) {
    return distance / motorCycle.speed;  
}
\end{codebox}

می‌بینید که مقدار زیادی کد تکراری داریم و مجبور شدیم برای هر وسیلهٔ نقلیه،
یک متد جدا بنویسیم. شرایط زمانی بدتر می‌شود که بخواهیم کلاس تریلی را هم
اضافه کنیم، اما یادمان برود که متد محاسبه زمان بر حسب مسافت را برای آن
بنویسیم.

ارث‌بری کار ما را خیلی راحت می‌کند چون زمانی که همهٔ کلاس‌ها از
\lr{\texttt{vehicle}} ارث‌بری کنند، می‌توانیم فقط یک متد برای محاسبهٔ زمان
داشته باشیم:

\begin{codebox}{java}
public static double calculateTime(double distance, Vehicle vehicle) {
    return distance / vehicle.speed;  
}
\end{codebox}

و این متد، توسط تمام آبجکت‌هایی که از جنس \lr{\texttt{Vehicle}} هستند،
قابل استفاده‌ است. این موضوع شامل آبجکت‌های \lr{\texttt{Car}} ،
\lr{\texttt{Bike}} ، \lr{\texttt{MotorCycle}} و حتی تریلی‌های آینده هم
می‌شود!

\begin{codebox}{java}
public class Main {
    public static void main(String[] args) {  
        Vehicle myVehicle = new Vehicle("Benz", 200);  
        Car myCar = new Car("Iran-Khodro", 180, 4);  
        Bike myBike = new Bike("Scott", 10, false);  
        MotorCycle myMotorCycle = new MotorCycle("Yamaha", 100, false);  
  
        System.out.println("time taken by myVehicle to move 10 km: " + calculateTime(10, myVehicle));  
        System.out.println("time taken by myCar to move 10 km: " + calculateTime(10, myCar));  
        System.out.println("time taken by myBike to move 10 km: " + calculateTime(10, myBike));  
        System.out.println("time taken by myMotorCycle to move 10 km: " + calculateTime(10, myMotorCycle));  
    }

    public static double calculateTime(double distance, Vehicle vehicle) {
        return distance / vehicle.speed;  
    }  
}
\end{codebox}

این کد، خروجی زیر را می‌دهد:

\begin{codebox}{text}
time taken by myVehicle to move 10 km: 0.05  
time taken by myCar to move 10 km: 0.05555555555555555  
time taken by myBike to move 10 km: 1.0  
time taken by myMotorCycle to move 10 km: 0.1
\end{codebox}

\subsubsection{\lr{explicit cast} در \lr{reference type} ها}

اگر دقت کنید، در \lr{cast} های بالا، لازم نبود که داخل پرانتز،
\lr{explicitly} مشخص کنیم تا جاوا بداند ما از مشکلات \lr{cast} آگاهیم،
چون \lr{cast} هایی که انجام دادیم اصلاً مشکلی به وجود نمی‌آوردند! ولی
همیشه این‌طور نیست. مثلاً اگر بخواهیم که یک \lr{\texttt{Vehicle}} را به یک
\lr{\texttt{Car}} تبدیل کنیم چه می‌شود؟

کد زیر را ببینید:

\begin{codebox}{java}
public static void main(String[] args) {  
    Car myCar = new Car("Iran-Khodro", 180, 4);  
    myCar.showCarInfo();  
      
    System.out.println();  
      
    Vehicle myVehicle = myCar;  
    myVehicle.showInfo();  
      
    System.out.println();  
      
    Car theSameCar = myVehicle;  
    theSameCar.showCarInfo();  
}
\end{codebox}

تلاش کنید آن را اجرا کنید. به همان خطای همیشگی می‌خورید:

\begin{figure}[H]
\centering
\includegraphics[width=0.8\textwidth]{images/oop-2-2.png}
\end{figure}

چرا شما باید بتونید \lr{\texttt{Vehicle}} رو به \lr{\texttt{Car}} تبدیل
کنید؟! ممکن است \lr{\texttt{Car}} فیلدهای اضافه نسبت به
\lr{\texttt{Vehicle}} داشته باشد. اما راه‌هایی برای انجام این \lr{cast}
وجود دارد. برای این که بفهمیم چرا و چطور باید نگاهی به مموری بکنیم. اگر
از بخش «reference \lr{type} ها»ی داک قبلی به خاطر داشته باشید، بعد از
درست شدن متغیر \lr{\texttt{myCar}} و مقدار‌دهی‌اش، مموری چنین شکلی دارد:

\begin{figure}[H]
\centering
\includegraphics[width=0.8\textwidth]{images/oop-2-3.png}
\end{figure}

بعد از این که \lr{\texttt{myVehicle}} هم ساخته شود، مموری به شکل زیر در
می‌آید:

\begin{figure}[H]
\centering
\includegraphics[width=0.8\textwidth]{images/oop-2-4.png}
\end{figure}

می‌توانید این موضوع را با تکه کد زیر تایید کنید. این‌جا می‌بینید که تغییر
سرعت \lr{\texttt{myVehicle}} ، باعث تغییر سرعت \lr{\texttt{myCar}} هم
می‌شود:

\begin{codebox}{java}
myVehicle.accelerate(100);  
  
System.out.println("myVehicle speed: " + myVehicle.speed);  
System.out.println("myCar speed: " + myCar.speed);
\end{codebox}

خروجی این کد به شکل زیر است:

\begin{codebox}{text}
myVehicle speed: 280  
myCar speed: 280
\end{codebox}

می‌بینید که سرعت هر دو دستگاه از \lr{\texttt{180}} به \lr{\texttt{280}}
رسیده. حالا سراغ \lr{cast} کردن می‌رویم:

\begin{codebox}{java}
Car theSameCar = myVehicle;
\end{codebox}

با این که جاوا اخطار مشکل‌زا بودن این \lr{cast} را می‌دهد، ولی اگر باز هم
به مموری نگاه کنیم، می‌بینیم که \lr{\texttt{myVehicle}} هم به آبجکتی از
جنس \lr{\texttt{Car}} اشاره می‌کند! آبجکت \lr{\texttt{car}} ای که
\lr{\texttt{myVehicle}} به آن اشاره می‌کند، همچنان فیلدی به اسم
\lr{\texttt{doors}} در مموری دارد و به همین خاطر، اگر به جاوا
\lr{explicitly} اطمینان دهیم که می‌دانیم چه کار می‌کنیم، این \lr{cast}
انجام می‌شود:

\begin{codebox}{java}
Car theSameCar = (Car)myVehicle;

theSameCar.showCarInfo();
\end{codebox}

اگر این کد را اجرا کنید، به خطایی نمی‌خورید و می‌بینید که
\lr{\texttt{theSameCar}} هم مثل \lr{\texttt{myCar}} چهار در دارد:

\begin{codebox}{text}
Brand: Iran-Khodro, Speed: 280 km/h  
Doors: 4
\end{codebox}

در نهایت مموری به شکل زیر خواهد بود:

\begin{figure}[H]
\centering
\includegraphics[width=0.8\textwidth]{images/oop-2-5.png}
\end{figure}

همان‌طور که می‌بینید، هر سه متغیری که در این برنامه تعریف کردید، دقیقاً به
یک‌جا اشاره می‌کنند.

\section{کلاس \lr{Object}}

\subsection{متدهای ناآشنا}

کد زیر را در \lr{\texttt{main}} بنویسید:

\begin{codebox}{java}
public static void main(String[] args) {  
    Car myCar = new Car("Iran-Khodro", 180, 4);  
    myCar.showCarInfo();  
  
    System.out.println("Hash code: " + myCar.hashCode());  
}
\end{codebox}

متد \lr{\texttt{hashCode()}} در خط آخر را کجا تعریف کردیم؟! چطور بدون
هیچ اخطار قبلی‌ای، از آن استفاده کردیم؟ خروجی این کد به شکل زیر است:

\begin{codebox}{text}
Brand: Iran-Khodro, Speed: 180 km/h  
Doors: 4  
Hash code: 500977346
\end{codebox}

همه چیز عجیب‌تر هم می‌شود! \lr{\texttt{myCar}} ، چندین متد دیگر هم دارد که
ما هیچ‌جا آن‌ها را تعریف نکردیم! حتی \lr{IntelliJ} هم به ما پیشنهاد می‌دهد
که از آن‌ها استفاده کنیم:

\begin{figure}[H]
\centering
\includegraphics[width=0.8\textwidth]{images/oop-2-6.png}
\end{figure}

روی همین متد \lr{\texttt{hashCode()}} برید و کلیک راست کنید، از منوی
\lr{Go To} ، گزینهٔ \lr{Declaration or Usages} را انتخاب کنید تا
پیاده‌سازی این متد را ببینید:

\begin{figure}[H]
\centering
\includegraphics[width=0.8\textwidth]{images/oop-2-7.png}
\end{figure}

اگر مقداری این صفحه را اسکرول کنید، می‌بینید که این متد، جزئی از کلاسی به
اسم \lr{Object} است:

\begin{figure}[H]
\centering
\includegraphics[width=0.8\textwidth]{images/oop-2-8.png}
\end{figure}

\subsection{تعریف کلاس \lr{Object}}

حالا کلاس \lr{\texttt{Object}} چی هست؟ \lr{\texttt{Object}} کلاسی‌ست که
همهٔ کلاس‌های دیگر از آن به صورت مستقیم یا غیرمستقیم ارث‌بری می‌کنند. در کد
شما، \lr{\texttt{Vehicle}} از \lr{\texttt{Object}} ارث‌بری می‌کند و
\lr{\texttt{Car}} هم از \lr{\texttt{Vehicle}} ، پس \lr{\texttt{Car}} به
شکلی غیرمستقیم از \lr{\texttt{Object}} ارث‌بری می‌کند و به همین خاطر، به
متدهای آن مثل \lr{\texttt{hashCode()}} ، \lr{\texttt{toString()}} یا
\lr{\texttt{equals(Object obj)}} دسترسی دارد (دقت کنید که کلاس
\lr{\texttt{Car}} همچنان فقط یک پدر دارد و کلاس \lr{\texttt{Object}} در
واقع پدربزرگ این کلاس است).

شما تا الآن بدون این که خبر داشته باشید از این متدها استفاده می‌کردید.
مثلا وقتی کد زیر را بزنید:

\begin{codebox}{java}
System.out.println(myCar);
\end{codebox}

خروجی زیر را می‌بینید:

\begin{codebox}{text}
Vehicle.Car@1ddc4ec2
\end{codebox}

این \lr{output} عجیب و غریب، در واقع خروجی متد
\lr{\texttt{myCar.toString()}} است. اگر کد زیر را اجرا کنید:

\begin{codebox}{java}
System.out.println(myCar.toString());
\end{codebox}

می‌بینید که خروجی هیچ تفاوتی ندارد:

\begin{codebox}{text}
Vehicle.Car@1ddc4ec2
\end{codebox}

حتی می‌توانید این موضوع را در پیاده‌سازی متد
\lr{\texttt{println(Object obj)}} هم ببینید. این متد به شکل زیر
پیاده‌سازی شده است:

\begin{codebox}{java}
public void println(Object x) {  
    String s = String.valueOf(x);
    if (getClass() == PrintStream.class) {
        // need to apply String.valueOf again since first invocation  
        // might return null
        writeln(String.valueOf(s));  
    } else {
        synchronized (this) {  
            print(s);  
            newLine();  
        }  
    }  
}
\end{codebox}

به جزئیات پیاده‌سازی این کد کاری نداریم فقط باید بدانیم که متد
\lr{\texttt{writeln}} در این کد \lr{output} را برای شما چاپ می‌کند.
می‌بینید که در این کد، \lr{\texttt{String.valueOf(x)}} صدا زده شده. این
متد هم پیاده‌سازی سادهٔ زیر را دارد:

\begin{codebox}{java}
public static String valueOf(Object obj) {
    return (obj == null) ? "null" : obj.toString();  
}
\end{codebox}

این متد، با ورودی گرفتن یک \lr{Object} ، اگر \lr{\texttt{null}} باشد،
رشتهٔ \lr{\texttt{null}} و در غیر این صورت \lr{\texttt{obj.toString()}}
را برمی‌گرداند.

این‌جا هم می‌توانید کمکی که \lr{inheritance} به ما می‌کند را ببینید. از
آن‌جایی که متد \lr{\texttt{println}} یک \lr{Object} ورودی می‌گیرد، شما
می‌توانید هر \lr{reference type} دیگری مثل \lr{\texttt{Vehicle}} ،
\lr{\texttt{Car}} و هر کلاس دیگری را به آن ورودی بدهید. فرض کنید مجبور
بودیم به خاطر عدم وجود \lr{inheritance} برای هر کلاس یک بار
\lr{\texttt{println}} را پیاده‌سازی کنیم!

\section{چه چیزی یاد گرفتیم؟}

مهم‌ترین نکاتی که یاد گرفتیم:

\begin{itemize}
\tightlist
\item
  این‌که \textbf{\lr{interface}} چیست و برای چه به آن نیاز داریم
\item
  مقایسۀ \textbf{\lr{interface} ها} و \textbf{\lr{abstract class} ها} و
  اینکه چرا نمی‌توان به‌صورت مستقیم از آنها شیء ساخت
\item
  خاصیت ارث‌بری چندگانه از اینترفیس‌ها
\item
  آشنایی با انواع \textbf{\lr{polymorphism} ها} و کاربرد آن‌ها
\item
  آشنایی با انواع و نحوۀ \textbf{\lr{casting}} در
  \textbf{\lr{primitive type} ها} و \textbf{\lr{refrence type} ها} و
  خطرات احتمالی آن
\item
  اتفاقات درون حافظه هنگام \textbf{\lr{casting}} در
  \textbf{\lr{refrences type} ها}
\item
  کلاس \textbf{\lr{object}} چیست
\item
  متدهای ناآشنایی که توسط ما تعریف نشده اند اما قابل استفاده اند، چه
  ارتباطی با کلاس \textbf{\lr{Object}} دارند
\end{itemize}
	
	
	
	
%%%%%%%%%%%%%%%%%%%%%%%%%%%%%%%%%%%%%%%%%%%%%%%%%%%%%%%%%%%%
\end{document}